\chapter{Introduction}

Urban mobility systems are increasingly undergoing a transition toward more sustainable and health-oriented transport solutions. Active travel modes such as walking and cycling are considered key components of this transition, as they contribute to reducing emissions while supporting more livable urban environments. 
Over the past two decades, cycling activity in Germany has risen substantially, with urban and educated populations contributing most strongly to this growth \cite{hudde_unequal_2022}. 
Research on sustainable mobility highlights that promoting active travel requires coordinated efforts in infrastructure, policy, and urban planning to influence travel behavior and support sustainable mobility patterns \cite{lattman_sustainable_2024}. 
Beyond environmental benefits, cycling also provides substantial public-health advantages. A systematic review of health impact assessments shows that increases in active transportation lead to considerable net health benefits, as gains from physical activity outweigh potential risks such as traffic incidents or air pollution exposure \cite{mueller_health_2015}. 
Consequently, urban and mobility planning increasingly promotes cycling as part of an integrated strategy that simultaneously addresses climate protection, accessibility, and public health \cite{kopal_healthy_2023}. As cycling becomes more prevalent in urban environments, understanding how cyclists interact within complex traffic systems becomes increasingly important.


Despite the benefits of cycling, urban environments continue to pose significant safety challenges for cyclists. Collision data from Germany show that densely populated areas experience higher rates of cyclist accidents, particularly at intersections and during turns, though these incidents tend to be less severe than those occurring in rural regions \cite{harkort_spatiotemporal_2023}. Fatal collisions, while less frequent in urban settings, often involve interactions with larger vehicles, highlighting the complexity of urban traffic risk \cite{bil_critical_2010}. Temporal factors further influence risk, with accident frequency peaking during morning rush hours on weekdays and varying across seasons.

Beyond individual risk, cyclist interactions themselves introduce additional safety considerations. Following and overtaking maneuvers, common in both solo and group riding, involve close proximity, speed variation, and rapid decision-making. Group cycling presents unique dynamics: cyclists benefit from increased visibility and shared situational awareness, but also face hazards due to reduced maneuvering space, social norms that may encourage risk-taking, and the collective influence of team behavior \cite{heeremans_group_2022}. These findings suggest that understanding not only individual behavior but also interaction patterns among cyclists is essential for effective safety interventions.


Cyclist behavior has been studied using a variety of complementary methods, each providing distinct perspectives on how riders interact with traffic and infrastructure. Eye-tracking studies typically examine visual attention, hazard perception, and cognitive workload in real or simulated environments, helping to understand how cyclists monitor their surroundings and respond to complex traffic situations \cite{kchour_understanding_2024, rupi_visual_2019, ma_eye_2024}. 
Bicycle simulators enable controlled investigations of riding dynamics and infrastructure interactions under repeatable conditions, supporting analysis of decision-making, lane positioning, and maneuver execution \cite{meuleners_improving_2023, sasaki_risks_2024, hammami_how_2025}. 
Complementing these approaches, instrumented bicycles capture high-resolution kinematic data in naturalistic or semi-controlled settings, allowing detailed assessment of speed, acceleration, lateral position, and responses to environmental factors \cite{belikhov_exploring_2025, mathuseck_bikesense_2024, saad_riderid_2025}. 
Surveys and perception studies, on the other hand, provide insight into cyclists’ subjective experiences, motivations, and risk perceptions, supporting the identification of different behavioral profiles and responses to infrastructure \cite{berghoefer_cyclists_2022, francke_are_2019}. 
While these methods have advanced understanding of cyclist behavior, they are often constrained by limited sample sizes, controlled or small-scale environments, or reliance on self-reported data.


Recent advances in aerial sensing have enabled large-scale, naturalistic observation of urban traffic, offering new opportunities to study cyclist behavior in context. Camera-equipped drones provide a bird’s-eye perspective over complex intersections, capturing interactions among vehicles, pedestrians, and cyclists without influencing participants’ behavior. Compared to ground-level sensors or controlled experiments, drones can monitor entire traffic areas continuously, reducing occlusion and providing temporally coherent trajectories across multiple intersections. This unobtrusive approach allows researchers to capture authentic motion patterns under real-world conditions.

Several datasets exemplify the potential of drone-based trajectory collection. The TUMDOT–MUC dataset \cite{kutsch_tumdotmuc_2024} tracked all road users along a 700m urban corridor in Munich for several hours across multiple drone viewpoints, capturing detailed kinematics such as position, velocity, acceleration, and orientation. Similarly, the inD dataset \cite{bock_ind_2020} records over 11,500 road users, including bicycles, across four urban intersections in Germany, enabling large-scale analysis of naturalistic interaction patterns. More recently, the HetroD dataset \cite{chen_hetrod_2026} emphasizes heterogeneous traffic dominated by vulnerable road users, providing more than 65,000 trajectories with high-fidelity annotations for modeling complex interactions. Beyond dataset creation, methodological work by \citeauthor{wang_method_2023} (\citeyear{wang_method_2023}) has demonstrated automated pipelines for extracting accurate trajectories from drone videos, combining detection, tracking, and cleaning steps to produce precise, continuous motion data suitable for downstream behavioral analysis.

Together, these resources enable microscopic investigations of traffic interactions at scales previously unattainable. By capturing continuous, high-resolution trajectories across multiple traffic modes, drone-based datasets allow detailed analysis of cyclist interactions in urban traffic. In particular, they provide access to fine-grained kinematic variables that describe short-term riding behavior. Following and overtaking maneuvers involve rapid adjustments in speed, spacing, and positioning, making the identification of local behavioral states an important step in understanding interaction dynamics.

Microscopic analysis of cyclist behavior examines short-term motion dynamics that form the foundation of maneuver execution. Key kinematic measures include longitudinal and lateral speed and acceleration, as well as yaw, capturing both forward control and directional changes \cite{arvin_how_2019, arvin_safety_2021, bachir_identifying_2025, li_volatility-based_2025, wang_what_2015, yarlagadda_exploring_2022}. These variables can be directly derived from trajectory data and provide detailed insight into individual rider responses to traffic and environmental conditions.

Unsupervised approaches have been used to segment trajectories into local behavioral states. \citeauthor{mohammed_characterization_2019} (\citeyear{mohammed_characterization_2019}) clustered following interactions into constrained and unconstrained states, and overtaking into initiation, merging, and post-overtaking states, analyzing multivariate distributions and state transitions. Trajectory segmentation has been implemented using event-based triggers, fixed or sliding windows, and context-aware dynamic windows \cite{yarlagadda_exploring_2022, bachir_identifying_2025, li_volatility-based_2025, zhang_proactive_2023}, allowing meaningful characterization of micro-level riding behavior.

Trajectory data has also been used to investigate specific interaction maneuvers among cyclists, particularly following and overtaking. These maneuvers involve continuous adjustments in spacing, speed, and positioning, making them suitable for analysis through high-resolution trajectory observations.

Following interactions among cyclists have been investigated using both empirical trajectory analysis and microscopic interaction modeling. Drone-based trajectory data have been used to analyze kinematic characteristics such as speed, acceleration, spacing, and headways during following and free-flow riding \cite{kutsch_analyzing_2025}. Beyond descriptive measurements, modeling approaches have been proposed to characterize spacing control and interaction dynamics. For example, stochastic headway models have been developed to represent the distribution of minimum and free headways maintained by cyclists \cite{hoogendoorn_bicycle_2016}, while unified follow-the-leader frameworks based on experimental trajectory data have been proposed to describe spacing–velocity relationships across multiple traffic modes \cite{zhao_unified_2017}. These approaches provide empirical measurements and theoretical models for microscopic following behavior.

Overtaking behavior has been examined through empirical trajectory observations, qualitative analyses, and automated detection methods. Trajectory-based studies have quantified kinematic characteristics of overtaking maneuvers, including speed differences, lateral spacing, and maneuver duration \cite{kutsch_analyzing_2025, khan_characteristics_2001}. Qualitative observational studies have also been used to describe overtaking strategies and inform traffic microsimulation models \cite{mocsari_analysis_2009}. More recently, algorithmic approaches have been developed to automatically detect overtaking events from trajectory data, for example through multi-phase detection frameworks applied to drone-based observations \cite{feng_automatic_2023}. In addition, statistical models have been used to analyze factors influencing overtaking decisions in naturalistic riding data \cite{wang_investigating_2022}.


Cyclist interactions in urban traffic, such as overtaking and following, are complex and shaped by both individual behavior and the surrounding environment. Despite increasing availability of drone-based trajectory data, systematic approaches to uncover recurring maneuver execution patterns remain limited. This thesis develops a multi-stage hierarchical unsupervised framework that segments maneuvers, captures internal kinematic dynamics, and aggregates them into maneuver-level representations, enabling the identification of distinct execution styles. Applied to overtaking and following maneuvers, the framework reveals characteristic temporal and kinematic patterns, bridging micro-level riding dynamics with maneuver-level behavioral structure while accounting for scenario exposure.